\documentclass[11pt]{article}
\usepackage[margin=1in]{geometry}
\usepackage{hyperref}
\usepackage{enumitem}
\title{Event-Driven Robot Coding}
\author{Piter Garcia}
\date{\today}
\begin{document}
\maketitle

\section*{Grade and Class Match}
\textbf{Grade:} Grade 4\\
\textbf{Likely blocks:} Day 1 · 1:15-2:10 · Baris; Day 2 · 1:15-2:10 · Fowler; Day 3 · 1:15-2:10 · Schrank; Day 4 · 1:15-2:10 · Autore

\section*{Lesson Focus}
Robotics with events, sensors, reactions, and challenge-based troubleshooting.

\section*{Learning Targets}
\begin{itemize}[leftmargin=*]
\item Students build an event-driven or sensor-based program.
\item Students test and revise code based on robot behavior.
\item Students explain the trigger and response in their program.
\end{itemize}

\section*{Materials}
\begin{itemize}[leftmargin=*]
\item Robots
\item Student devices
\item Challenge setup or obstacle space
\end{itemize}

\section*{Suggested Lesson Flow}
\begin{enumerate}[leftmargin=*]
\item Open with the exact device, tool, or routine students need in order to start successfully.
\item Model one example task before students begin independent or partner work.
\item Give students time to build, code, design, or document their work while circulating for support.
\item Close with a short share-out, reflection, or debugging discussion.
\end{enumerate}

\section*{Source Notes}
This lesson packet is enriched with transcript-derived lesson notes, the school schedule roster, and official starting-point resources. See the paired Markdown guide for clickable links.

\bibliographystyle{plain}
\bibliography{event-driven-robot-coding}
\end{document}
